% 表:失效覆盖矩阵,2026-07-21 改版。行是 3.1 节五类失效,列是八个机制,列头旋转、分 framework 与 benchmark 两组,
% 勾表示"该机制收窄该失效"。图案即对账:每行至少一个勾;噪声行三个勾;机制名与 3.2/3.3 正文一致。
\begin{table}[t]
\scriptsize
\centering
\caption{Coverage of the five failure sources by the mechanisms of the framework and the benchmark.}
% 表:framework 与 benchmark 的机制对五类失效源的覆盖。
\label{tab:coverage}
\setlength{\tabcolsep}{3.5pt}
\resizebox{\columnwidth}{!}{%
\begin{tabular}{lcccccccc}
\hline
 & \multicolumn{3}{c}{Framework} & \multicolumn{5}{c}{Benchmark} \\
Failure source
 & \rotatebox{90}{Entry points}
 & \rotatebox{90}{Measurements}
 & \rotatebox{90}{Disp.\ kernels}
 & \rotatebox{90}{Verifier}
 & \rotatebox{90}{Hidden checks}
 & \rotatebox{90}{Control band}
 & \rotatebox{90}{Threshold}
 & \rotatebox{90}{Replay} \\
\hline
Fabricated reports & & \checkmark & & & & & & \\
Tampering with the environment & \checkmark & & & & & & & \\
Wrong semantics & & & & \checkmark & \checkmark & & & \\
Crashing the validator & & & \checkmark & & & & & \\
Noise reported as gains & & & & & & \checkmark & \checkmark & \checkmark \\
\hline
\end{tabular}%
}
\end{table}
